
\begin{refsection}%[chapters/Ref/intro.bib]
\refstepcounter{subsubsection}
\chapter{Introduction, literature review, research gap and objective}


\begin{figure}[H]
	\centering
	\includegraphics[width=0.8\textwidth]{chapters/graphs/Introduction/Picture1.png}
	\caption{Different pathways involved in immunosuppressive microenvironment through MDSCs recruitment and TAMs secreting TGF-$\upbeta$, CCL-22, PD-L1, and IL-10 to promote T-regulatory function and arginine depletion. Tumor immune evasion through PD-1/PD-L1 and CTLA-4 checkpoint activation, IDO-mediated tryptophan catabolism, and adenosine pathway signaling via CD39/CD73, collectively suppressing CTL and NK cell activity. These mechanisms converge to downregulate MHC-I expression, reduce DC antigen presentation, and inhibit cytotoxic T lymphocyte activation, differentiation, and proliferation.}
	\label{fig:i1}
\end{figure}



\begin{table}[htbp]
	\centering
	\caption{Different types of Quantum dots, their origin, synthesis and biomedical applications}
	\label{sheet:QD}
	\small
	\renewcommand{\arraystretch}{1.15}
	\begin{tabularx}{\textwidth}{|c|>{\raggedright\arraybackslash}X|>{\raggedright\arraybackslash}X|>{\raggedright\arraybackslash}X|>{\raggedright\arraybackslash}X|p{0.8cm}|}
		\hline
		\textbf{S.No.} & \textbf{Type} & \textbf{Origin} & \textbf{Synthesis} & \textbf{Biomedical applications} & \textbf{Ref.} \\
		\hline
		1. & Carbon quantum dots (CQDs) & CQDs derived from carbon precursors; often grouped with diamond dots & Hydrothermal, microwave, electro chemical, laser, and pyrolysis methods & Bioimaging, biosensing, drug delivery, and theranostics & 1 \\
		\hline
		2. & Graphene quantum dots (GQDs) & Graphene derived carbon nanomaterials & Top-down cutting of graphene/ graphite materials or bottom-up synthesis using hydrothermal, microwave, and electro chemical routes & Bioimaging, drug delivery, biosensing, tissue engineering, antibacterial/ antiviral uses, and theranostics & 1 \\
		\hline
		3. & InP QDs & Indium phosphide based semiconductor QDs & Synthesized mainly by colloidal methods and core/shell growth with surface passivation & Biomedical fluorescence imaging, immunoassays, biosensing, drug delivery, and diagnostics & 1 \\
		\hline
		4. & Cd-based QDs & CdSe and CdTe & Organo metallic, water-based, or microbial based bottom up synthesis; core/shell designs to improve stability and reduce toxicity & Imaging, diagnostics, and other biomedical uses & 1\\
		\hline
		5. & Cd-free quantum dots & InP based, silicon, and other cadmium free systems & Diverse synthesis methods & Biophotonics, nanomedicine, bioimaging, sensing, drug/gene delivery, and theranostics &1\\
		\hline
		6. & Silicon quantum dots (SiQDs) & Silicon & Laser ablation and other synthesis approaches & In vivo imaging, drug delivery, and tissue engineering & 1\\
		\hline
	\end{tabularx}
\end{table}





\section{Main objectives of this dissertation}

On the basis of the above discussion about the review of literature, related definition of the problem, contemporary challenges, and drawbacks of existing therapies along with the need for new improvement.
Following main objectives of this dissertation are elucidated below in brief.
\begin{itemize}
	\item \textbf{Objective 1:} Objective 1.
	\item \textbf{Objective 2:} Objective 2.
	\item \textbf{Objective 3:} Objective 3.
\end{itemize}

\textbf{Organization of this dissertation:} This dissertation comprises the following sequence of chapters:
\begin{enumerate}[label=\textbf{Chapter \arabic*:}]
	\item Introduction and review of literature.
	
	\item Materials, methods, equipment design, and computational approaches.
	
	\item Biological experiment results, discussion of obtained results, and major findings.
	
	\item In Chapter 4 of this dissertation, we have summarized the results obtained through this dissertation are
	concluded. The future scope of this work also has been included at the end of this chapter.
\end{enumerate}

\printbibliography
\end{refsection}